\chapter{Analysis}\label{ch:analysis}

In this chapter, we describe in detail functional and non-functional requirements of a solution for the problem.

\section{Functional requirements}
% Which functions must be offered to users / other programs?  Which are the input data and the output data? Which is the expected effect?
No specific interface is required: the system should be able to operate without any external interface. In order to conduct tests and evaluate the system, a minimal service for performing link and node operations is provided, as well as a monitoring service for observing the system's state.


\section{Non functional requirements}

The system must be fault tolerant, in particular it must have the property of eventual consistency: edge and node can fail or be temporarily unavailable, so the service must keep operating and converge to a correct state after a finite amount of time. The whole system should be transparent to the user: they should be unaware, or at least not be bothered by, the elimination or addition of edges during communications based on the MST.

% TODO: CAP theorem


% Cheng et al. : PA:
% - eventual consistency: the system reaches a consistent state after a finite time
% - availability (even if possibly not optimal)
% - partition tolerance: non-negotiable since we want a dynamic mst

% Everything about mode and transparencies: availability, mobility, security, fault tolerance, etc.



% Are there execution time bounds? Minimum data rates?

% If requested, specific platforms/languages/middlewares requirements for the implementation can be decided here. (E.g.: if the project is on a SOA, we may request that functions are offered via SOAP or RESTful services).
